\phantomsection\addcontentsline{toc}{chapter}{课程资料与教学进度}
\noindent{\color{structurecolor}\bfseries\fontsize{14.4}{20}\selectfont 课程资料与教学进度}\par
主课程为 MIT 数学系本科课程18.703“现代代数”, James McKernan 教授, 2013年春季学期. 官方安排每周两次讲课, 每次1.5小时. 官网先修列为18.02多元微积分; 18.700线性代数与18.703共同组成标准代数序列. 本中文册在模和域部分使用线性代数基础, 这是补充来源带来的阅读要求, 不改写成主课程官方先修.

大纲为有限群理论, 环、理想、欧几里得环和多项式的唯一分解, 域及有限域的性质和应用. 指定教材为 I. N. Herstein, \textit{Abstract Algebra}, Macmillan, 1986, ISBN 9780023538209. 阅读页另列 Tom Judson 的在线教材 \textit{Abstract Algebra: Theory and Applications} 各节作为课前阅读. 仅有书目题号的作业不代表官网公开了相应教材题面.

考核为作业50\%, 两次期中各10\%, 期末30\%. 课程说明作业通常周二布置、下一周周四交, 考试周可例外; 不收迟交, 去掉最低一次作业分. 考试有复述定义部分. 以下按照官方教学进度表的讲次和 E1/E2/E3 考试标签列出, 官网未列具体日历日期. 原作业文件的截止日期另随题集登记, 不用它反推每次讲课日期.
\begin{longtable}{@{}p{.09\textwidth}p{.45\textwidth}p{.16\textwidth}p{.22\textwidth}@{}}
\toprule 讲次&主题&交作业&Judson阅读节\\\midrule\endhead
1&群&&3.1,3.2\\
2&子群&&3.3\\
3&陪集&&6.1,6.2\\
4&循环群&作业1&6.2,4.1\\
5&置换群&&4.2,5.1\\
6&对称群中的共轭&&5.1,5.2\\
7&同构&作业2&9.1\\
8&同态与核&&11.1,10.1\\
9&商群&作业3&10.1,11.2\\
10&同构定理&&11.2\\
11&交错群&作业4&9.2\\
12&表示与小阶群&&13.1\\
13&Sylow定理与应用&作业5&15.1,15.2\\
14&环&作业6&16.1,16.2\\
E1&第一次期中&&\\
15&环的基本性质&&期中复习 (阅读页原列)\\
16&环同态与理想&作业7&16.3\\
17&分式域&作业8&18.1,16.4\\
18&素理想与极大理想&&18.2\\
19&特殊整环&作业9&18.2,17.2\\
20&欧几里得整环&&17.2,17.3\\
E2&第二次期中&作业10&\\
21&多项式环&&17.2,17.3\\
22&快速素性检验&作业11&17.3\\
23&群作用与自同构&&期中复习 (阅读页原列)\\
24&复习, 无讲义&&17.3\\
E3&期末考试&&\\
\bottomrule
\end{longtable}
补充资料 RES.18-011 是2021年秋季18.701“代数I”的学生笔记, 讲课教师 Davesh Maulik, 记录者 Jakin Ng、Sanjana Das、Ethan Yang, 合订178页. RES.18-012 是2022年春季18.702“代数II”的学生笔记, 讲课教师 Roman Bezrukavnikov, 记录者 Sanjana Das、Jakin Ng, 笔记工作由 Ashay Athalye 监督, 合订162页. 两个资源页均明确未经教师核准; 此处“监督记录”不等于审定数学内容. 两册各35讲, 第二册另有特征标附录. 没有在这两组资源的页面公开独立作业或考试卷; 不把其所链接的其他年份课程材料冒充2021/2022年的试题. 本册只按明确选定的补充讲次使用.

参考: \href{https://ocw.mit.edu/courses/18-703-modern-algebra-spring-2013/pages/syllabus/}{官方课程大纲}, \href{https://ocw.mit.edu/courses/18-703-modern-algebra-spring-2013/pages/calendar/}{教学进度}, \href{https://ocw.mit.edu/courses/18-703-modern-algebra-spring-2013/pages/readings/}{阅读安排}.
